% Lösungen Einheit 2 von 5 – Winkel an Geradenkreuzungen und Parallelen (zusammenbau v0.1)
\einheitenkopf{Einheit 2 von 5 · Winkel an Geradenkreuzungen und Parallelen}

\erg{12}{a) als parallel angegeben}
%% TODO Lösung der Erklärzeile 13g) fehlt in der Bank
\erg{13}{a) Scheitelwinkel \quad b) $38^\circ$ (Scheitelwinkel sind gleich groß) \quad c) $51^\circ$ (Scheitelwinkel sind gleich groß) \quad d) $29^\circ$ (Scheitelwinkel sind gleich groß) \quad e) $66^\circ$ (Scheitelwinkel sind gleich groß) \quad f) $83^\circ$ (Scheitelwinkel sind gleich groß) \quad g) \ldots \quad h) $180^\circ - 44^\circ = 136^\circ$ \quad i) $\beta = 108^\circ$, $\gamma = 72^\circ$, $\delta = 108^\circ$ \quad j) Scheitelwinkel $64^\circ$; $\beta = 64^\circ - 23^\circ = 41^\circ$ \quad k) $58^\circ$ (Stufenwinkel) \quad l) $49^\circ$ (Wechselwinkel) \quad m) Stufenwinkel an h $62^\circ$; gesucht: $180^\circ - 62^\circ = 118^\circ$ \quad n) $\beta = 43^\circ$ (Scheitelwinkel); die Angabe $58^\circ$ wird nicht gebraucht \quad o) $\beta = 116^\circ$, $\gamma = 64^\circ$, $\delta = 116^\circ$ \quad p) Stufenwinkel an g $44^\circ$; $\alpha = 180^\circ - 44^\circ = 136^\circ$}
\erg{14}{a) ja; die Stufenwinkel sind gleich groß}
\erg{15}{a) Fehler: Nebenwinkel statt Scheitelwinkel berechnet; richtig: $\gamma = 66^\circ$ \quad b) α und γ sind beide Nebenwinkel von β: $\alpha = 180^\circ - \beta$ und $\gamma = 180^\circ - \beta$, also ist $\alpha = \gamma$. \quad c) $62^\circ$ (Stufenwinkel) und $180^\circ - 62^\circ = 118^\circ$}
